\documentclass[12pt, twoside, italian]{book}
%%%%%%%%%%%%%%%%%%%%%%%%%%%%%%%%%
% PACKAGE IMPORTS
%%%%%%%%%%%%%%%%%%%%%%%%%%%%%%%%%


\usepackage[tmargin=2cm,bmargin=2cm,footskip=.2in, rmargin=1.5in,lmargin=1.5in, margin=0.85in, a4paper]{geometry}
\usepackage{bookmark}
\usepackage{libertinus}
\usepackage{amsmath,amsfonts,amsthm,amssymb,mathtools}
\usepackage[varbb]{newpxmath}
%\usepackage[libertine]{newtxmath}
\usepackage{chemfig}
\usepackage{ragged2e}
\usepackage{xfrac}
\usepackage{mhchem}
\usepackage[italian]{babel}
\usepackage[Sonny]{fncychap}
\usepackage[makeroom]{cancel}
\usepackage{mathtools}
\usepackage{listing}
\usepackage{enumitem}
\usepackage{hyperref,theoremref}
\hypersetup{
	pdftitle={Appunti di Complementi di Analisi Matematica 2},
	colorlinks=true, linkcolor=doc!90,
	bookmarksnumbered=true,
	bookmarksopen=true,
	pdfauthor={Francesco Sermi}
}
\usepackage[most,many,breakable]{tcolorbox}
\usepackage{xcolor}
\usepackage{varwidth}
\usepackage{varwidth}
\usepackage{etoolbox}
%\usepackage{authblk}
\usepackage{nameref}
\usepackage{multicol,array}
\usepackage{tikz}
\usepackage{tikz-cd}
\usepackage{tikz-3dplot}
\usepackage{tikz}
\usetikzlibrary{3d,decorations.pathmorphing}
\usepackage[ruled,vlined,linesnumbered]{algorithm2e}
\usepackage{import}
\usepackage{xifthen}
\usepackage{pdfpages}
\usepackage{transparent}

\newcommand\mycommfont[1]{\footnotesize\ttfamily\textcolor{blue}{#1}}
\SetCommentSty{mycommfont}
\newcommand{\incfig}[1]{%
    \def\svgwidth{\columnwidth}
    \import{./figures/}{#1.pdf_tex}
}

\usepackage{tikzsymbols}
\usepackage{float}
\usepackage[toc, page]{appendix}
\renewcommand\qedsymbol{$\square$}
\usepackage{hyperref}
\usepackage{mwe}
\usepackage{fancyhdr}
\usepackage{pgfplots}
\usepackage{bm}
\usepackage{dsfont}
%\usepackage{import}
%\usepackage{xifthen}
%\usepackage{pdfpages}
%\usepackage{transparent}
\definecolor{doc}{RGB}{0,0,0}
\theoremstyle{definition}
\newtheorem{definition}{Definizione}[chapter]

\theoremstyle{definition}
\newtheorem{theorem}{Teorema}[chapter]
\theoremstyle{definition}
\newtheorem{cor}{Corollario}[theorem]
\theoremstyle{definition}
\newtheorem{lemma}{Lemma}[theorem]
\theoremstyle{plain}
\newtheorem{exercise}{Esercizio}[chapter]
\theoremstyle{definition}
\newtheorem{prop}{Proposizione}[chapter]

\theoremstyle{definition}
\newtheorem{example}{Esempio}[chapter]

\theoremstyle{remark}
\newtheorem*{remark}{Oss:}

\newcommand*\circled[1]{\tikz[baseline=(char.base)]{
		\node[shape=circle,draw,inner sep=1pt] (char) {#1};}}
\newcommand{\innerprod}[2]{\langle #1, #2 \rangle}
\renewcommand\appendixpagename{Appendici}
\renewcommand\appendixtocname{Appendici}
\setcounter{localmathalphabets}{0}
% Things Lie
\newcommand{\kb}{\mathfrak b}
\newcommand{\kg}{\mathfrak g}
\newcommand{\kh}{\mathfrak h}
\newcommand{\kn}{\mathfrak n}
\newcommand{\ku}{\mathfrak u}
\newcommand{\kz}{\mathfrak z}
\DeclareMathOperator{\Ext}{Ext} % Ext functor
\DeclareMathOperator{\Tor}{Tor} % Tor functor
\newcommand{\gl}{\opname{\mathfrak{gl}}} % frak gl group
\renewcommand{\sl}{\opname{\mathfrak{sl}}} % frak sl group chktex 6

% More script letters etc.
\newcommand{\SA}{\mathcal A}
\newcommand{\SB}{\mathcal B}
\newcommand{\SC}{\mathcal C}
\newcommand{\SF}{\mathcal F}
\newcommand{\SG}{\mathcal G}
\newcommand{\SH}{\mathcal H}
\newcommand{\OO}{\mathcal O}

\newcommand{\SCA}{\mathscr A}
\newcommand{\SCB}{\mathscr B}
\newcommand{\SCC}{\mathscr C}
\newcommand{\SCD}{\mathscr D}
\newcommand{\SCE}{\mathscr E}
\newcommand{\SCF}{\mathscr F}
\newcommand{\SCG}{\mathscr G}
\newcommand{\SCH}{\mathscr H}

% Mathfrak primes
\newcommand{\km}{\mathfrak m}
\newcommand{\kp}{\mathfrak p}
\newcommand{\kq}{\mathfrak q}

% number sets
\newcommand{\RR}[1][]{\ensuremath{\ifstrempty{#1}{\mathbb{R}}{\mathbb{R}^{#1}}}}
\newcommand{\NN}[1][]{\ensuremath{\ifstrempty{#1}{\mathbb{N}}{\mathbb{N}^{#1}}}}
\newcommand{\ZZ}[1][]{\ensuremath{\ifstrempty{#1}{\mathbb{Z}}{\mathbb{Z}^{#1}}}}
\newcommand{\QQ}[1][]{\ensuremath{\ifstrempty{#1}{\mathbb{Q}}{\mathbb{Q}^{#1}}}}
\newcommand{\CC}[1][]{\ensuremath{\ifstrempty{#1}{\mathbb{C}}{\mathbb{C}^{#1}}}}
\newcommand{\PP}[1][]{\ensuremath{\ifstrempty{#1}{\mathbb{P}}{\mathbb{P}^{#1}}}}
\newcommand{\HH}[1][]{\ensuremath{\ifstrempty{#1}{\mathbb{H}}{\mathbb{H}^{#1}}}}
\newcommand{\FF}[1][]{\ensuremath{\ifstrempty{#1}{\mathbb{F}}{\mathbb{F}^{#1}}}}
% expected value
\newcommand{\EE}{\ensuremath{\mathbb{E}}}
\newcommand{\charin}{\text{ char }}
\DeclareMathOperator{\sign}{sign}
\DeclareMathOperator{\Aut}{Aut}
\DeclareMathOperator{\Inn}{Inn}
\DeclareMathOperator{\Syl}{Syl}
\DeclareMathOperator{\Gal}{Gal}
\DeclareMathOperator{\GL}{GL} % General linear group
\DeclareMathOperator{\SL}{SL} % Special linear group

%---------------------------------------
% BlackBoard Math Fonts :-
%---------------------------------------

%Captital Letters
\newcommand{\bbA}{\mathbb{A}}	\newcommand{\bbB}{\mathbb{B}}
\newcommand{\bbC}{\mathbb{C}}	\newcommand{\bbD}{\mathbb{D}}
\newcommand{\bbE}{\mathbb{E}}	\newcommand{\bbF}{\mathbb{F}}
\newcommand{\bbG}{\mathbb{G}}	\newcommand{\bbH}{\mathbb{H}}
\newcommand{\bbI}{\mathbb{I}}	\newcommand{\bbJ}{\mathbb{J}}
\newcommand{\bbK}{\mathbb{K}}	\newcommand{\bbL}{\mathbb{L}}
\newcommand{\bbM}{\mathbb{M}}	\newcommand{\bbN}{\mathbb{N}}
\newcommand{\bbO}{\mathbb{O}}	\newcommand{\bbP}{\mathbb{P}}
\newcommand{\bbQ}{\mathbb{Q}}	\newcommand{\bbR}{\mathbb{R}}
\newcommand{\bbS}{\mathbb{S}}	\newcommand{\bbT}{\mathbb{T}}
\newcommand{\bbU}{\mathbb{U}}	\newcommand{\bbV}{\mathbb{V}}
\newcommand{\bbW}{\mathbb{W}}	\newcommand{\bbX}{\mathbb{X}}
\newcommand{\bbY}{\mathbb{Y}}	\newcommand{\bbZ}{\mathbb{Z}}

%---------------------------------------
% MathCal Fonts :-
%---------------------------------------

%Captital Letters
\newcommand{\mcA}{\mathcal{A}}	\newcommand{\mcB}{\mathcal{B}}
\newcommand{\mcC}{\mathcal{C}}	\newcommand{\mcD}{\mathcal{D}}
\newcommand{\mcE}{\mathcal{E}}	\newcommand{\mcF}{\mathcal{F}}
\newcommand{\mcG}{\mathcal{G}}	\newcommand{\mcH}{\mathcal{H}}
\newcommand{\mcI}{\mathcal{I}}	\newcommand{\mcJ}{\mathcal{J}}
\newcommand{\mcK}{\mathcal{K}}	\newcommand{\mcL}{\mathcal{L}}
\newcommand{\mcM}{\mathcal{M}}	\newcommand{\mcN}{\mathcal{N}}
\newcommand{\mcO}{\mathcal{O}}	\newcommand{\mcP}{\mathcal{P}}
\newcommand{\mcQ}{\mathcal{Q}}	\newcommand{\mcR}{\mathcal{R}}
\newcommand{\mcS}{\mathcal{S}}	\newcommand{\mcT}{\mathcal{T}}
\newcommand{\mcU}{\mathcal{U}}	\newcommand{\mcV}{\mathcal{V}}
\newcommand{\mcW}{\mathcal{W}}	\newcommand{\mcX}{\mathcal{X}}
\newcommand{\mcY}{\mathcal{Y}}	\newcommand{\mcZ}{\mathcal{Z}}


%---------------------------------------
% Bold Math Fonts :-
%---------------------------------------

%Captital Letters
\newcommand{\bmA}{\boldsymbol{A}}	\newcommand{\bmB}{\boldsymbol{B}}
\newcommand{\bmC}{\boldsymbol{C}}	\newcommand{\bmD}{\boldsymbol{D}}
\newcommand{\bmE}{\boldsymbol{E}}	\newcommand{\bmF}{\boldsymbol{F}}
\newcommand{\bmG}{\boldsymbol{G}}	\newcommand{\bmH}{\boldsymbol{H}}
\newcommand{\bmI}{\boldsymbol{I}}	\newcommand{\bmJ}{\boldsymbol{J}}
\newcommand{\bmK}{\boldsymbol{K}}	\newcommand{\bmL}{\boldsymbol{L}}
\newcommand{\bmM}{\boldsymbol{M}}	\newcommand{\bmN}{\boldsymbol{N}}
\newcommand{\bmO}{\boldsymbol{O}}	\newcommand{\bmP}{\boldsymbol{P}}
\newcommand{\bmQ}{\boldsymbol{Q}}	\newcommand{\bmR}{\boldsymbol{R}}
\newcommand{\bmS}{\boldsymbol{S}}	\newcommand{\bmT}{\boldsymbol{T}}
\newcommand{\bmU}{\boldsymbol{U}}	\newcommand{\bmV}{\boldsymbol{V}}
\newcommand{\bmW}{\boldsymbol{W}}	\newcommand{\bmX}{\boldsymbol{X}}
\newcommand{\bmY}{\boldsymbol{Y}}	\newcommand{\bmZ}{\boldsymbol{Z}}
%Small Letters
\newcommand{\bma}{\boldsymbol{a}}	\newcommand{\bmb}{\boldsymbol{b}}
\newcommand{\bmc}{\boldsymbol{c}}	\newcommand{\bmd}{\boldsymbol{d}}
\newcommand{\bme}{\boldsymbol{e}}	\newcommand{\bmf}{\boldsymbol{f}}
\newcommand{\bmg}{\boldsymbol{g}}	\newcommand{\bmh}{\boldsymbol{h}}
\newcommand{\bmi}{\boldsymbol{i}}	\newcommand{\bmj}{\boldsymbol{j}}
\newcommand{\bmk}{\boldsymbol{k}}	\newcommand{\bml}{\boldsymbol{l}}
\newcommand{\bmm}{\boldsymbol{m}}	\newcommand{\bmn}{\boldsymbol{n}}
\newcommand{\bmo}{\boldsymbol{o}}	\newcommand{\bmp}{\boldsymbol{p}}
\newcommand{\bmq}{\boldsymbol{q}}	\newcommand{\bmr}{\boldsymbol{r}}
\newcommand{\bms}{\boldsymbol{s}}	\newcommand{\bmt}{\boldsymbol{t}}
\newcommand{\bmu}{\boldsymbol{u}}	\newcommand{\bmv}{\boldsymbol{v}}
\newcommand{\bmw}{\boldsymbol{w}}	\newcommand{\bmx}{\boldsymbol{x}}
\newcommand{\bmy}{\boldsymbol{y}}	\newcommand{\bmz}{\boldsymbol{z}}

%---------------------------------------
% Scr Math Fonts :-
%---------------------------------------

\newcommand{\sA}{{\mathscr{A}}}   \newcommand{\sB}{{\mathscr{B}}}
\newcommand{\sC}{{\mathscr{C}}}   \newcommand{\sD}{{\mathscr{D}}}
\newcommand{\sE}{{\mathscr{E}}}   \newcommand{\sF}{{\mathscr{F}}}
\newcommand{\sG}{{\mathscr{G}}}   \newcommand{\sH}{{\mathscr{H}}}
\newcommand{\sI}{{\mathscr{I}}}   \newcommand{\sJ}{{\mathscr{J}}}
\newcommand{\sK}{{\mathscr{K}}}   \newcommand{\sL}{{\mathscr{L}}}
\newcommand{\sM}{{\mathscr{M}}}   \newcommand{\sN}{{\mathscr{N}}}
\newcommand{\sO}{{\mathscr{O}}}   \newcommand{\sP}{{\mathscr{P}}}
\newcommand{\sQ}{{\mathscr{Q}}}   \newcommand{\sR}{{\mathscr{R}}}
\newcommand{\sS}{{\mathscr{S}}}   \newcommand{\sT}{{\mathscr{T}}}
\newcommand{\sU}{{\mathscr{U}}}   \newcommand{\sV}{{\mathscr{V}}}
\newcommand{\sW}{{\mathscr{W}}}   \newcommand{\sX}{{\mathscr{X}}}
\newcommand{\sY}{{\mathscr{Y}}}   \newcommand{\sZ}{{\mathscr{Z}}}


%---------------------------------------
% Math Fraktur Font
%---------------------------------------

%Captital Letters
\newcommand{\mfA}{\mathfrak{A}}	\newcommand{\mfB}{\mathfrak{B}}
\newcommand{\mfC}{\mathfrak{C}}	\newcommand{\mfD}{\mathfrak{D}}
\newcommand{\mfE}{\mathfrak{E}}	\newcommand{\mfF}{\mathfrak{F}}
\newcommand{\mfG}{\mathfrak{G}}	\newcommand{\mfH}{\mathfrak{H}}
\newcommand{\mfI}{\mathfrak{I}}	\newcommand{\mfJ}{\mathfrak{J}}
\newcommand{\mfK}{\mathfrak{K}}	\newcommand{\mfL}{\mathfrak{L}}
\newcommand{\mfM}{\mathfrak{M}}	\newcommand{\mfN}{\mathfrak{N}}
\newcommand{\mfO}{\mathfrak{O}}	\newcommand{\mfP}{\mathfrak{P}}
\newcommand{\mfQ}{\mathfrak{Q}}	\newcommand{\mfR}{\mathfrak{R}}
\newcommand{\mfS}{\mathfrak{S}}	\newcommand{\mfT}{\mathfrak{T}}
\newcommand{\mfU}{\mathfrak{U}}	\newcommand{\mfV}{\mathfrak{V}}
\newcommand{\mfW}{\mathfrak{W}}	\newcommand{\mfX}{\mathfrak{X}}
\newcommand{\mfY}{\mathfrak{Y}}	\newcommand{\mfZ}{\mathfrak{Z}}
%Small Letters
\newcommand{\mfa}{\mathfrak{a}}	\newcommand{\mfb}{\mathfrak{b}}
\newcommand{\mfc}{\mathfrak{c}}	\newcommand{\mfd}{\mathfrak{d}}
\newcommand{\mfe}{\mathfrak{e}}	\newcommand{\mff}{\mathfrak{f}}
\newcommand{\mfg}{\mathfrak{g}}	\newcommand{\mfh}{\mathfrak{h}}
\newcommand{\mfi}{\mathfrak{i}}	\newcommand{\mfj}{\mathfrak{j}}
\newcommand{\mfk}{\mathfrak{k}}	\newcommand{\mfl}{\mathfrak{l}}
\newcommand{\mfm}{\mathfrak{m}}	\newcommand{\mfn}{\mathfrak{n}}
\newcommand{\mfo}{\mathfrak{o}}	\newcommand{\mfp}{\mathfrak{p}}
\newcommand{\mfq}{\mathfrak{q}}	\newcommand{\mfr}{\mathfrak{r}}
\newcommand{\mfs}{\mathfrak{s}}	\newcommand{\mft}{\mathfrak{t}}
\newcommand{\mfu}{\mathfrak{u}}	\newcommand{\mfv}{\mathfrak{v}}
\newcommand{\mfw}{\mathfrak{w}}	\newcommand{\mfx}{\mathfrak{x}}
\newcommand{\mfy}{\mathfrak{y}}	\newcommand{\mfz}{\mathfrak{z}}

%From M275 "Topology" at SJSU
\newcommand{\id}{\mathrm{id}}
\newcommand{\taking}[1]{\xrightarrow{#1}}
\newcommand{\inv}{^{-1}}

%From M170 "Introduction to Graph Theory" at SJSU
\DeclareMathOperator{\diam}{diam}
\DeclareMathOperator{\ord}{ord}
\newcommand{\defeq}{\overset{\mathrm{def}}{=}}

%From the USAMO .tex files
\newcommand{\ts}{\textsuperscript}
\newcommand{\dg}{^\circ}
\newcommand{\ii}{\item}

% % From Math 55 and Math 145 at Harvard
% \newenvironment{subproof}[1][Proof]{%
% \begin{proof}[#1] \renewcommand{\qedsymbol}{$\blacksquare$}}%
% {\end{proof}}

\newcommand{\liff}{\leftrightarrow}
\newcommand{\lthen}{\rightarrow}
\newcommand{\opname}{\operatorname}
\newcommand{\surjto}{\twoheadrightarrow}
\newcommand{\injto}{\hookrightarrow}
\newcommand{\On}{\mathrm{On}} % ordinals
\DeclareMathOperator{\img}{im} % Image
\DeclareMathOperator{\Img}{Im} % Image
\DeclareMathOperator{\coker}{coker} % Cokernel
\DeclareMathOperator{\Coker}{Coker} % Cokernel
\DeclareMathOperator{\Ker}{Ker} % Kernel
\DeclareMathOperator{\rank}{rank}
\DeclareMathOperator{\Spec}{Spec} % spectrum
\DeclareMathOperator{\Tr}{Tr} % trace
\DeclareMathOperator{\pr}{pr} % projection
\DeclareMathOperator{\ext}{ext} % extension
\DeclareMathOperator{\pred}{pred} % predecessor
\DeclareMathOperator{\dom}{dom} % domain
\DeclareMathOperator{\ran}{ran} % range
\DeclareMathOperator{\Hom}{Hom} % homomorphism
\DeclareMathOperator{\Mor}{Mor} % morphisms
\DeclareMathOperator{\End}{End} % endomorphism

\newcommand{\eps}{\epsilon}
\newcommand{\veps}{\varepsilon}
\newcommand{\ol}{\overline}
\newcommand{\ul}{\underline}
\newcommand{\wt}{\widetilde}
\newcommand{\wh}{\widehat}
\newcommand{\vocab}[1]{\textbf{\color{blue} #1}}
\providecommand{\half}{\frac{1}{2}}
\newcommand{\dang}{\measuredangle} %% Directed angle
\newcommand{\ray}[1]{\overrightarrow{#1}}
\newcommand{\seg}[1]{\overline{#1}}
\newcommand{\arc}[1]{\wideparen{#1}}
\DeclareMathOperator{\cis}{cis}
\DeclareMathOperator*{\lcm}{lcm}
\DeclareMathOperator*{\argmin}{arg min}
\DeclareMathOperator*{\argmax}{arg max}
\newcommand{\cycsum}{\sum_{\mathrm{cyc}}}
\newcommand{\symsum}{\sum_{\mathrm{sym}}}
\newcommand{\cycprod}{\prod_{\mathrm{cyc}}}
\newcommand{\symprod}{\prod_{\mathrm{sym}}}
\newcommand{\Qed}{\begin{flushright}\qed\end{flushright}}
\newcommand{\parinn}{\setlength{\parindent}{1cm}}
\newcommand{\parinf}{\setlength{\parindent}{0cm}}
% \newcommand{\norm}{\|\cdot\|}
\newcommand{\inorm}{\norm_{\infty}}
\newcommand{\opensets}{\{V_{\alpha}\}_{\alpha\in I}}
\newcommand{\oset}{V_{\alpha}}
\newcommand{\opset}[1]{V_{\alpha_{#1}}}
\newcommand{\lub}{\text{lub}}
\newcommand{\del}[2]{\frac{\partial #1}{\partial #2}}
\newcommand{\Del}[3]{\frac{\partial^{#1} #2}{\partial^{#1} #3}}
\newcommand{\deld}[2]{\dfrac{\partial #1}{\partial #2}}
\newcommand{\Deld}[3]{\dfrac{\partial^{#1} #2}{\partial^{#1} #3}}
\newcommand{\lm}{\lambda}
\newcommand{\uin}{\mathbin{\rotatebox[origin=c]{90}{$\in$}}}
\newcommand{\usubset}{\mathbin{\rotatebox[origin=c]{90}{$\subset$}}}
\newcommand{\lt}{\left}
\newcommand{\rt}{\right}
\newcommand{\bs}[1]{\boldsymbol{#1}}
\newcommand{\exs}{\exists}
\newcommand{\st}{\strut}
\newcommand{\dps}[1]{\displaystyle{#1}}

\newcommand{\sol}{\setlength{\parindent}{0cm}\textbf{\textit{Soluzione:}}\setlength{\parindent}{1cm} }
\newcommand{\solve}[1]{\setlength{\parindent}{0cm}\textbf{\textit{Soluzione: }}\setlength{\parindent}{1cm}#1 \Qed}
\DeclareMathOperator{\Div}{div}
\DeclareMathOperator{\Rot}{rot}
\newcommand{\dcdot}{%
  \mathbin{%
    \nonscript\mspace{-\muexpr\medmuskip*2/3}%
    \cdot
    \nonscript\mspace{-\muexpr\medmuskip*2/3}%
  }%
}

\usepackage{fancyhdr}
\usepackage{titlesec}  % Opzionale, per controllo avanzato sui titoli
\setlength{\parindent}{0pt}
\setlength{\parskip}{0pt}
\pagestyle{fancy}

% Definizione dell'intestazione per pagine dispari (Odd)
\fancyhead[LO]{\textbf{\thepage}}  % Numero di pagina in grassetto a sinistra
\fancyhead[RO]{\nouppercase\rightmark}         % Capitolo e sezione a destra

% Definizione dell'intestazione per pagine pari (Even)
\fancyhead[LE]{\nouppercase\leftmark}          % Capitolo e sezione a sinistra
\fancyhead[RE]{\textbf{\thepage}}  % Numero di pagina in grassetto a destra

% Cancella il footer
\fancyfoot{}
\setlength{\headheight}{14.49998pt}
\addtolength{\topmargin}{-2.49998pt}
% Personalizzazione della linea orizzontale
\renewcommand{\headrulewidth}{0.4pt}
\renewcommand{\footrulewidth}{0pt}

\pgfplotsset{compat=newest}
\title{Appunti di Geometria e Algebra Lineare}
\author{Francesco Sermi}
\date{\today}
\pgfplotsset{
	compat = newest,
	colormap/violet % se si vuole cambiare il colormap utilizzato per generare i grafici, tuttavia la copertina ha un colormap a
					% parte che va modificato qua sotto. Se si stampa non a colori consiglio blackwhite come colormap
	}
\renewcommand{\appendixpage}{}
\begin{document}
	\begin{titlepage}
	\centering
	\vspace*{3cm}
	{\huge\bfseries Appunti di Geometria e Algebra Lineare \par}
	\vspace{2cm}
	{\Large\itshape Francesco Sermi\par}
	\vfill
	\begin{figure}[H]
		\centering
		%\includegraphics[scale=0.125]{immagini/grey_burningship.png}
	\end{figure}
    	%\begin{tikzpicture}
	% Definisco lo stile dei vettori
    	%\tikzset{vector/.style={-stealth, thick, color=black}}

    % Parametri per il punto di contatto
    	%\def\xo{1}
    	%\def\yo{2}
    	%\def\zo{5} % Calcola z0 = f(x0, y0)

    % Derivate parziali della funzione f(x, y) = x^2 + y^2 nel punto (x0, y0)
    	%\def\fx{2*\xo} % = 2
    	%\def\fy{2*\yo} % = 4

    % Disegno la superficie
    	%\begin{axis}[
    	%    width=12cm,
    	%    view={60}{30}, % Angolazione impostata
    	%    axis lines=middle,
    	%    zlabel={$f(x,y)$},
    	%    clip=false,
    	%    domain=-5:5,
    	%    y domain=-5:5,
    	%    samples=30,
    	%    colormap/viridis,
    	%    z buffer=sort
    	%]

    % Grafico della superficie con colore visibile
    	%\addplot3[surf, opacity=0.8, shader=interp]
    	%    {x^2 - y^2}; % Superficie f(x, y) = x^2 + y^2
    	%\end{axis}
	%\end{tikzpicture}
	\vfill	
	\end{titlepage}
	\pagestyle{empty}
	\vspace*{\fill}
	Da ficcarci il copyright.
	\cleardoublepage
	\vspace*{8cm}
	\begin{center}
			\begin{flushright}
				\small\emph{Al mio amico e compagno di banco Giulio}
			\end{flushright}
		\vspace*{\fill}
	\end{center}
	\cleardoublepage
	\pagestyle{fancy}
	\tableofcontents
	\chapter*{Introduzione}
	\pagestyle{empty}
	\thispagestyle{empty}
	\addcontentsline{toc}{chapter}{Introduzione}
	\markboth{Introduzione}{Introduzione}
	\pagestyle{fancy}
	Queste note nascono come appunti di Geometria e Algebra Lineare	
	\begin{center}
		francesco291104 (\textbf{dot}) gmail.com
	\end{center}

	oppure fermandomi direttamente in dipartimento. In ogni caso, vi auguro di appassionarvi allo studio di questa materia.\

	Buono studio a tutti! \\ \\ \\ \\
	\emph{Pisa, \today} \hfill Francesco Sermi\
	\vfill
	\cleardoublepage
	\chapter*{Notazioni}
	\markboth{Notazioni}{Notazioni}
	\pagestyle{empty}
	\thispagestyle{empty}
	\pagestyle{fancy}

	Riporto all'inizio del libro le notazioni adottate all'interno di questo documento:
	\begin{itemize}[label=\hspace{-0.5em}]
		\item $\underline{v}$ vettore
		\item $||\cdot||, |\cdot|$ norma
		\item $\innerprod{\cdot}{\cdot}$ prodotto scalare/hermitiano
		\item $a \times b, a \wedge b$ prodotto vettoriale
		\item $|| A ||$ norma
		\item $\sum_{n}$ gruppo (rispetto all'operazione di composizione di funzioni) delle permutazioni $\sigma: \{1, \ldots n \} \to \{1, \ldots n \}$ bigettive.
		\item $\iota_+$ segnatura positiva di un prodotto scalare/hermitiano
		\item $\iota_-$ segnatura negativa di un prodotto scalare/hermitiano
		\item $\iota_0$ segnatura nulla di un prodotto scalare/hermitiano
		\item $(\iota_+, \iota_-, \iota_0)$ segnatura di un prodotto scalare/hermitiano
		\item $I(x_0, \delta) = (x_0 - \delta, x_0 + \delta) = B(x_0, \delta) \subset \mathbb{R}$ con $x_0 \in \mathbb{R}, \delta > 0$ intervallo centrato in $x_0$ di ampiezza $2\delta$
		\item $\varprod\limits_i$ prodotto cartesiano multiplo
		\item $\sum\limits_i$ sommatoria
		\item $\sim$ relazione d'equivalenza
		\item $\sfrac{X}{\sim}$ insieme quoziente di $X$ rispetto alla relazione d'equivalenza
 
	\end{itemize}
	
	\newpage
	\chapter{Spazi vettoriali}

\thispagestyle{empty}

Il cuore del corso di Geometria e Algebra Lineare è sicuramente lo studio degli spazi vettoriali. In questa sezione vedremo le definizioni, le proprietà e i teoriemi fondamentali che riguardano gli spazi vettoriali, i sottospazi, il concetto di base e di dimensione. \\

\section{Definizioni e proprietà}
La definizione di spazio vettoriale è strettamente legata alla definizione di campo, che può essere così definito
\begin{definition}[Campo]
    \label{def:campo}
    Un \textbf{campo} è un insieme $\mathbb{K}$ dotato di due operazioni $+_\mathbb{K}: \mathbb{K} \times \mathbb{K} \to \mathbb{K}$ e $\cdot_\mathbb{K}: \mathbb{K} \times \mathbb{K} \to \mathbb{K}$ tali che:
    \begin{itemize}
        \item $(\mathbb{K}, +)$ è un gruppo abeliano, ovvero l'operazione $+$ gode delle seguenti proprietà:
        \begin{enumerate}[label=\protect\circled{\arabic*}]
            \item $\forall a, b, c \in \mathbb{K}, (a +_\mathbb{K} b) +_\mathbb{K} c = a +_\mathbb{K} (b +_\mathbb{K} c)$ (proprietà associativa);
            \item $\forall a, b \in \mathbb{K}, a +_\mathbb{K} b = b +_\mathbb{K} a$ (proprietà commutativa);
            \item $\exists 0_\mathbb{K} \in \mathbb{K} : \forall a \in \mathbb{K}, 0_\mathbb{K} +_\mathbb{K} a = a +_\mathbb{K} 0_\mathbb{K} = a$ (esistenza dell'elemento neutro);
            \item $\forall a \in \mathbb{K}, \exists (-a) \in \mathbb{K} : a +_\mathbb{K} (-a) = (-a) +_\mathbb{K} a = 0$ (esistenza dell'elemento inverso).
        \end{enumerate}
        \item $(\mathbb{K} \setminus \{ 0_\mathbb{K} \}, \cdot)$ è un gruppo abeliano, ovvero l'operazione $\cdot$ gode delle seguenti proprietà:
        \begin{enumerate}[label=\protect\circled{\arabic*}]
            \item $\forall a, b, c \in \mathbb{K}, (a \cdot_\mathbb{K} b) \cdot_\mathbb{K} c = a \cdot_\mathbb{K} (b \cdot_\mathbb{K} c)$ (proprietà associativa);
            \item $\forall a, b \in \mathbb{K}, a \cdot_\mathbb{K} b = b \cdot_\mathbb{K} a$ (proprietà commutativa);
            \item $\exists \mathds{1}_{\mathbb{K}} \in \mathbb{K} : \mathds{1}_{\mathbb{K}} \cdot_\mathbb{K} a = a \cdot_\mathbb{K} \mathds{1}_{\mathbb{K}} = a$ (esistenza dell'elemento neutro);
            \item $\forall a \in \mathbb{K} \setminus \{ 0_\mathbb{K} \}, \exists a^{-1} \in \mathbb{K} : a \cdot_\mathbb{K} a^{-1} = a^{-1} \cdot_\mathbb{K} a = 1$ (esistenza dell'elemento inverso).
        \end{enumerate}
    \end{itemize}
\end{definition}
\begin{remark}
    Fare attenzione ai pedici sulle operazioni e sugli elementi: infatti, nel corso di questi appunti avremo a che fare spesso con delle operazioni che sono definite su insiemi diversi ma che vengono preservate dai morfismi; tuttavia, per evitare confusione, cercherò il più possibile di denotare le operazioni e gli elementi con pedici che tolgano ogni dubbio sull'insieme di appartenenza.
\end{remark}
\begin{example}
    L'esempio più lampante e semplice di campo è chiaramente l'insieme dei numeri razionali $\mathbb{Q}$ e l'insieme dei numeri reali $\mathbb{R}$ (e chiaramente anche il campo dei numeri complessi $\mathbb{C}$). Un esempio di campo meno noto è sicuramente $\mathbb{Z}_2$, ovvero l'insieme dei numeri interi modulo $2$, che contiene solo due elementi, $0$ e $1$, e le operazioni di somma e prodotto sono definite come segue:
    \begin{itemize}
        \item $0 + 0 = 0$, $0 + 1 = 1$, $1 + 0 = 1$, $1 + 1 = 0$;
        \item $0 \cdot 0 = 0$, $0 \cdot 1 = 0$, $1 \cdot 0 = 0$, $1 \cdot 1 = 1$.
    \end{itemize}
    e si può verificare che tale insieme è effettivamente un campo, in quanto le due operazioni possiedono tutte le proprietà richieste dalla definizione data in \ref{def:campo}.
\end{example}
\begin{definition}[Spazio vettoriale]
    \label{def:sp_vett} Sia dato un insieme $V$, un campo $\mathbb{K}$ e siano date due operazioni $+: V \times V \to V$ e $\cdot: \mathbb{K} \times V \to V$. Diremo che $(V, \mathbb{K}, +, \cdot)$ è un $\mathbb{K}-$spazio vettoriale (o più semplicemente uno spazio vettoriale sul campo $\mathbb{K}$) se le due operazioni soddisfano le seguenti proprietà:
    \begin{itemize}
        \item $(V, +)$ è un gruppo abeliano, ovvero l'operazione $+$ gode delle seguenti proprietà:
        \begin{enumerate}[label=\protect\circled{\arabic*}]
            \item $\forall \underline{u}, \underline{v}, \underline{z} \in V, (\underline{u} + \underline{v}) + \underline{z} = \underline{u} + (\underline{v} + \underline{z})$ (proprietà associativa);
            \item $\forall \underline{v}, \underline{w} \in V, \underline{v} + \underline{w} = \underline{w} + \underline{v}$ (proprietà commutativa);
            \item $\exists \underline{0} \in V : \forall \underline{v} \in V, \underline{0} + \underline{v} = \underline{v} + \underline{0} = \underline{v}$ (esistenza dell'elemento neutro);
            \item $\forall \underline{v} \in V, \exists (-\underline{v}) \in V : \underline{v} + (-\underline{v}) = (-\underline{v}) + \underline{v} = \underline{0}$ (esistenza dell'elemento inverso).
        \end{enumerate}
        \item l'operazione $\cdot: \mathbb{K} \times V \to V$ gode delle seguenti proprietà:
        \begin{enumerate}[label=\protect\circled{\arabic*}]
            \item $\forall \alpha \in \mathbb{K}, \forall \underline{v}, \underline{w} \in V, \alpha \cdot (\underline{v} + \underline{w}) = \alpha \cdot \underline{v} + \alpha \cdot \underline{w}$ (proprietà distributiva rispetto alla $+$);
            \item $\forall \alpha, \beta \in \mathbb{K}, \forall \underline{v} \in V, (\alpha +_{\mathbb{K}} \beta) \cdot \underline{v} = \alpha \cdot \underline{v} + \beta \cdot \underline{v}$ (proprietà distributiva rispetto alla $+_\mathbb{K}$);
            \item $\forall \alpha, \beta \in \mathbb{K}, \forall \underline{v} \in V, (\alpha \beta) \cdot \underline{v} = \alpha \cdot (\beta \underline{v})$ (proprietà associativa della $\cdot$);
            \item $\exists \mathds{1} \in \mathbb{K} : \forall \underline{v} \in V, \mathds{1} \cdot \underline{v} = \underline{v}$ (esistenza dell'elemento neutro per $\cdot$).
        \end{enumerate}
        Le due operazioni $+$ e $\cdot$ vengono dette, rispettivamente, \textbf{somma} (o legge di composizione interna) e \textbf{prodotto per scalare} (o legge di composizione esterna).
    \end{itemize}
\end{definition}
\begin{remark}
    Concettualmente, sebbene siano simili, la definizione di spazio vettoriale è diversa da quella di campo: infatti, sul campo abbiamo due operazioni che sono definite sul prodotto cartesiano del campo con sé stesso e restituiscono un elemento del campo; mentre nel caso dello spazio vettoriale si ha un insieme $V$ su cui sono definite due operazioni, una che è definita sul prodotto cartesiano di $V$ con sé stesso e restituisce un elemento di $V$, e l'altra che è definita sul prodotto cartesiano del campo con l'insieme $V$ e restituisce un elemento di $V$.
\end{remark}
\begin{remark}
    Si noti che uno spazio vettoriale è definito completamente solamente se definiamo le due operazioni $+$ e $\cdot$, oltre al campo. E' possibile, come vedremo in un esempio, andare a fornire la struttura di spazio vettoriale ad un medesimo insieme definendo le due operazioni in modo opportuno.
\end{remark}
Tramite questa definizione, possiamo andare a mostrare alcuni risultati validi nei campi che sono altrettanto validi negli spazi vettoriali.
\begin{prop}[l'elemento neutro rispetto a $+$ è unico]
    Sia $(V, \mathbb{K}, +, \cdot)$ uno spazio vettoriale sul campo $\mathbb{K}$. Allora l'elemento neutro rispetto all'operazione $+$ è unico, ovvero
    $$
    \exists! \underline{0} \in V : \forall \underline{v} \in V, \underline{0} + \underline{v} = \underline{v} + \underline{0} = \underline{v}.
    $$
\end{prop}
\begin{prop}
    \label{prop:0v_equal_0} Sia $(V, \mathbb{K}, +, \cdot)$ uno spazio vettoriale sul campo $\mathbb{K}$. Allora
    $$
    \forall \underline{v} \in V, 0_{\mathbb{K}} \cdot \underline{v} = \underline{0}.
    $$
\end{prop}
\begin{prop}
    Sia $(V, \mathbb{K}, +, \cdot)$ uno spazio vettoriale sul campo $\mathbb{K}$. Allora
    $$
    \forall \underline{v} \in V, (-1_\mathbb{K}) \cdot \underline{v} = (-\underline{v}).
    $$
\end{prop}
\begin{exercise}
    \label{ex:1.1} Mostrare le tre proposizioni precedenti. (\emph{Suggerimento}: per la prima si può procedere per assurdo, negando la tesi; per la seconda si può mostrare direttamente usando il fatto che $0_\mathbb{K} = 0_\mathbb{K} +_\mathbb{K} 0_\mathbb{K}$ e la terza si può mostrare direttamente usando la seconda.)
\end{exercise}
Vediamo adesso qualche esempio di spazio vettoriale per fissare meglio i concetti.
\begin{example}[$\mathbb{R}^2$ è uno spazio vettoriale sul campo $\mathbb{R}$]
    Consideriamo l'insieme $\mathbb{R}^2 = \mathbb{R} \times \mathbb{R}$: gli elementi di questo insieme sono coppie ordinate di numeri reali, che possiamo denotare come $\underline{v} = (v_1, v_2)$ con $v_1, v_2 \in \mathbb{R}$. Possiamo definire le due operazioni $+$ e $\cdot$ come segue:
    \begin{align*}
        &\underline{v} + \underline{w} = \begin{pmatrix} v_1 \\ v_2 \end{pmatrix} + \begin{pmatrix} w_1 \\ w_2 \end{pmatrix} = \begin{pmatrix} v_1 + w_1 \\ v_2 + w_2 \end{pmatrix}
        &\alpha \cdot \underline{v} = \alpha \cdot \begin{pmatrix} v_1 \\ v_2 \end{pmatrix} = \begin{pmatrix} \alpha v_1 \\ \alpha v_2 \end{pmatrix}
    \end{align*}
    ed è possibile verificare che queste due operazioni soddisfano sempre le proprietà richieste dalla definizione \ref{def:sp_vett}
\end{example}
\begin{example}[$\mathbb{R}^n$ è uno spazio vettoriale sul campo $\mathbb{R}$]
    Partendo dall'idea dell'esempio precedente, possiamo iterare il ragionamento e definire l'insieme $\mathbb{R}^n = \underbrace{\mathbb{R} \times \ldots \times \mathbb{R}}_{\text{n volte}}$: gli elementi di questo insieme sono $n$-tuple ordinate di numeri reali, che possiamo denotare come $\underline{v} = (v_1, \ldots, v_n)$ con $v_i \in \mathbb{R} \forall i \in \{ 1, \ldots, n \}$. Possiamo definire la somma e il prodotto per scalare come fatto prima, sommando componente per componente e moltiplicando componente per componente:
    \begin{align*}
        &\underline{v} + \underline{w} = \begin{pmatrix} v_1 \\ \vdots \\ v_n \end{pmatrix} + \begin{pmatrix} w_1 \\ \vdots \\ w_n \end{pmatrix} = \begin{pmatrix} v_1 + w_1 \\ \vdots \\ v_n + w_n \end{pmatrix} \\
        &\alpha \cdot \underline{v} = \alpha \cdot \begin{pmatrix} v_1 \\ \vdots \\ v_n \end{pmatrix} = \begin{pmatrix} \alpha v_1 \\ \vdots \\ \alpha v_n \end{pmatrix}.
    \end{align*}
\end{example}
\begin{example}[$\mathbb{C}^n$ è uno spazio vettoriale sul campo $\mathbb{C}$]
    Analogamente a quanto fatto per $\mathbb{R}^n$, possiamo riadattare il ragionamento fatto per $\mathbb{R}^n$ ammettendo che le componenti delle $n$-tuple ordinate siano numeri complessi, ovvero $\mathbb{C}^n \ni \underline{v} = \begin{pmatrix} v_1 \\ \vdots \\ v_n \end{pmatrix}$ con $v_i \in \mathbb{C} \forall i \in \{ 1, \ldots, n \}$.
\end{example}
\begin{example}[le soluzioni di un'equazione differenziale lineare sono uno spazio vettoriale]
    Sia data un'equazione differenziale lineare di ordine $n$ del tipo
    $$
        \sum_i^n a_i y^{(i)}(x) = 0 \text{ con } a_i \in \mathbb{R}
    $$
    allora possiamo mostrare che l'insieme delle soluzioni di questa equazione, denotato con $V$ gode della struttura di spazio vettoriale sul campo $\mathbb{R}$ definendo le due operazioni come segue: prese due soluzioni $y_1, y_2 \in V$ allora poniamo
    \begin{align*}
        &(y_1 + y_2)(x) = y_1(x) + y_2(x) \, \, \forall x \in \mathbb{R}
        &(\alpha y_1)(x) = \alpha y_1(x) \, \, \forall x \in \mathbb{R}
    \end{align*}
    e osserviamo che queste due operazioni sono chiuse, ovvero restituiscono un elemento che è ancora una soluzione dell'equazione differenziale lineare. Infatti
    $$
        \sum_i^n a_i (y_1 + y_2)^{(i)}(x) = \sum_i^n a_i y_1^{(i)} + \sum_i^n a_i y_2^{(i)} = 0 + 0 = 0,
    $$
    dove abbiamo usato il fatto che l'operazione di derivata è distributiva rispetto alla somma e il fatto che $y_1, y_2$ fossero soluzioni dell'equazione differenziale lineare, ovvero $\sum_i^n a_i y_j^{(i)} = 0$, con $j=1, 2$. Similmente, per quanto riguarda il prodotto per scalare, si usa il fatto che le costanti non sono coinvolte nel processo di derivazione, dunque
    $$
        \sum_i^n a_i(\alpha y_1)^{(i)}(x) = \sum_i^n a_i \alpha y_1^{(i)}(x) = \alpha \sum_i^n a_i y_1^{(i)}(x) = 0,
    $$
    dove si è sempre usato il fatto che $y_1$ fosse una soluzione dell'equazione differenziale lineare. Le proprietà richieste dalla definizione \ref{def:sp_vett} sono banalmente verificabili.
\end{example}
\begin{example}[spazio dei polinomi a coefficienti reali di grado $\leq n$ è uno spazio vettoriale]
    Sia dato l'insieme dei polinomi a coefficienti reali con grado $\leq n$ $\mathbb{R}_{\leq n}[x]$ su $\mathbb{R}$, allora possiamo definire le due operazioni come segue: presi due polinomi $p_1, p_2 \in \mathbb{R}_{\leq n}[x]$ possiamo definire il seguente polinomio
    \begin{align*}
    &(p_1 + p_2)(x) = p_1(x) + p_2(x) \, \, \forall x \in \mathbb{R}
    &(\alpha p_1)(x) = \alpha p_1(x) \, \, \forall x \in \mathbb{R}.
    \end{align*}
    Si può verificare che queste due operazioni sono chiuse, ovvero restituiscono sempre un polinomio a coefficienti reali di grado $\leq n$: infatti se $p_1(x) = \sum\limits_{i=1}^n a_i x^i$ e $p_2(x) = \sum\limits_{i=1}^n b_i x^i$ con $a_i, b_i \in \mathbb{R}$ allora
    \begin{align*}
        &(p_1 + p_2)(x) = \sum_{i=1}^n a_i x^i + \sum_{i=1}^n b_i x^i = \sum_{i=1}^n (a_i + b_i) x^i \\
        &(\alpha p_1)(x) = \alpha \sum_{i=1}^n a_i x^i = \sum_{i=1}^n (\alpha a_i) x^i
    \end{align*}
\end{example}
\begin{example}[Uno spazio vettoriale è definito solamente conoscendo la legge interna ed esterna]
    Fissiamo come spazio vettoriale $V = \mathbb{R}^n$, che è chiaramente uno spazio vettoriale sul campo $\mathbb{R}$, e definiamo le due operazioni $\oplus$ e $\odot$ come
    \begin{align*}
        &\underline{v} \oplus \underline{w} = \begin{pmatrix} v_1 \\ \vdots \\ v_n \end{pmatrix} \oplus \begin{pmatrix} w_1 \\ \vdots \\ w_n \end{pmatrix} = \begin{pmatrix} v_1 + w_1 + 2 \\ \vdots \\ v_n + w_n +2 \end{pmatrix} \\
        &\alpha \odot \underline{v} = \alpha \odot \begin{pmatrix} v_1, \ldots, v_n \end{pmatrix} = \begin{pmatrix} \alpha v_1 - 2(1-\alpha) \\ \alpha v_2 - 2(1-\alpha) \\ \vdots \\ \alpha v_n - 2(1-\alpha) \end{pmatrix}.
    \end{align*}
    Si può verificare che queste due operazioni sono chiuse, ovvero restituiscono sempre un elemento di $V$, e soddisfano le proprietà richieste dalla definizione \ref{def:sp_vett} di spazio vettoriale. Ne deduciamo che uno spazio vettoriale è completamente determinato \textbf{solo} conoscendo le due operazioni $+: V \times V \to V$ e $\cdot: \mathbb{K} \times V \to V$, sebbene in questo spazio vettoriale non è il vettore $(0, \ldots, 0)$ ad essere l'elemento neutro rispetto all'operazione $+$, ma il vettore $(-2, \ldots, -2)$; tuttavia, è interessante notare come
    comunque valgano le proprietà dimostrate nella proposizione \ref{prop:0v_equal_0} (infatti, $0 \odot \underline{v} = \begin{pmatrix} 0 v_1 - 2(1-0) \\ \vdots \\ 0v_n - 2(1-0) \end{pmatrix} = \begin{pmatrix} -2 \\ \vdots \\ -2 \end{pmatrix} )$ e nelle due successivi mostrate.
\end{example}
\section{Sottospazi}

Siamo ora interessati a studiare dei sottoinsiemi di uno spazio vettoriale che siano a loro volta degli spazi vettoriali rispetto all'operazione di somma e di prodotto per scalare. I sottospazi più comuni che possiamo immaginare, come vedremo, sono le rette e i piani che abbiamo in $\mathbb{R}^3$ (le rette anche in $\mathbb{R}^2$): la potenza dell'approccio che svilupperemo è il fatto che potremo trattare oggetti, come gli iperpiani, come se fossero oggetti i nostri usuali piani. 
\begin{definition}[sottospazio]
    Sia $(V, \mathbb{K}, +, \cdot)$ uno spazio vettoriale sul campo $\mathbb{K}$ e sia $W \subseteq V$ un sottoinsieme di $V$. Diremo che $W$ è un \textbf{sottospazio vettoriale} se esso è uno spazio vettoriale sul campo $\mathbb{K}$ rispetto alle operazioni $+$ e $\cdot$ definite su $V$.
\end{definition}
\begin{remark}
    Ogni spazio vettoriale è un sottospazio di sé stesso, siccome è chiaro che $V \subseteq V$.
\end{remark}
\begin{prop}
    \label{prop:ssp_vett_iff}
    Sia $(V, \mathbb{K}, +, \cdot)$ uno spazio vettoriale sul campo $\mathbb{K}$. Allora $W \subseteq V$ è un sottospazio vettoriale se e solo se
    \begin{enumerate}[label=\protect\circled{\arabic*}]
        \item $\forall \underline{v}, \underline{w} \in V, (\underline{v}, \underline{w} \in W) \implies \underline{v} + \underline{w} \in W$;
        \item $\forall \alpha \in \mathbb{K}, \forall \underline{v} \in V, (\underline{v} \in W) \implies (\alpha \underline{v} \in W)$.
    \end{enumerate}
\end{prop}
\begin{proof} Mostriamo le due implicazioni. \\
    $\boxed{\Rightarrow}$: se $W \subseteq V$ è un sottospazio vettoriale, allora per definizione esso è uno spazio vettoriale sul campo $\mathbb{K}$, dunque le due implicazioni
    \begin{align*}
        &\underline{v}, \underline{w} \in W \implies \underline{v} + \underline{w} \in W;
        &\underline{v} \in W \implies \alpha \underline{v} \in W;
    \end{align*}
    sono sempre vere. \\
    $\boxed{\Leftarrow}$: se valgono le due implicazioni, allora è chiaro che le due operazioni $+$ e $\cdot$ sono chiuse su $W$, dunque è sufficiente mostrare che valgono le proprietà richiesta dalla definizione \ref{def:sp_vett} di spazio vettoriale. Tuttavia, siccome queste valgono per ogni elemento di $V$, è chiaro che varranno anche per ogni elemento di $W$ siccome $W \subseteq V$; dunque possiamo concludere che $W$ è un sottospazio vettoriale di $V$.  
\end{proof}
\textbf{Attenzione}: da adesso in poi, per comodità di scrittura, non scriverò più esplicitamente l'operazione $\cdot$ per indicare il prodotto per scalare, ma mi limiterò a scrivere direttamente $\alpha \underline{v}$ sottointendendo il prodotto per scalare.
\begin{definition}
    Sia $(V, \mathbb{K}, +, \cdot)$ uno spazio vettoriale sul campo $\mathbb{K}$ e siano $\underline{v_1}, \ldots, \underline{v_n} \in V$, definiamo l'insieme delle combinazioni lineari di $\underline{v_1}, \ldots, \underline{v_n}$ come
    $$
        \text{Span}(\underline{v_1}, \ldots, \underline{v_n}) = \left\{ \underline{v} \in V \, \bigg| \, \exists \alpha_1, \ldots, \alpha_n \in \mathbb{K} : \underline{v} = \alpha_1 \underline{v_1} + \ldots + \alpha_n \underline{v_n} \right\}
    $$
\end{definition}
\begin{prop}
    Sia $(V, \mathbb{K}, +, \cdot)$ uno spazio vettoriale sul campo $\mathbb{K}$ e siano $\underline{v_1}, \ldots, \underline{v_n} \in V$. Allora $\text{Span}(\underline{v_1}, \ldots, \underline{v_n})$ è un sottospazio vettoriale di $V$.
\end{prop}
\begin{prop}
    Sia $(V, \mathbb{K}, +, \cdot)$ uno spazio vettoriale sul campo $\mathbb{K}$ e siano $\underline{v_1}, \ldots, \underline{v_n} \in V$. Allora
    $$
        \text{Span}(\underline{v_1}, \ldots, \underline{v_n}) \text{ è il più piccolo sottospazio vettoriale di } V \text{ che contiene } \underline{v_1}, \ldots, \underline{v_n}.
    $$
\end{prop}
\begin{exercise}
    \label{ex:1.2} Mostrare le due proposizioni qua sopra. 
\end{exercise}
Possiamo generalizzare la definizione di $\text{Span}$ di vettori ad insiemi di vettori arbitrari, ovvero
\begin{definition}[Span di un insieme]
    Sia $(V, \mathbb{K}, +, \cdot)$ uno spazio vettoriale sul campo $\mathbb{K}$ e sia $A \subseteq V$ un insieme, allora definiamo lo \text{Span} di $A$ come
    $$
        \text{Span}(A) = \left\{ \underline{v} \in V \, \bigg| \, \exists \underline{v_1}, \ldots, \underline{v_n} \in A : \exists \alpha_1, \ldots, \alpha_n \in \mathbb{K} : \underline{v} = \alpha_1 \underline{v_1} + \ldots \alpha_n \underline{v_n} \right\}.
    $$
\end{definition}
\begin{prop}[Span di un insieme è un sottospazio vettoriale]
    Sia $(V, \mathbb{K}, +,\cdot)$ uno spazio vettoriale sul campo $\mathbb{K}$ e sia $A \subseteq V$ un insieme. Allora $\text{Span}(A)$ è un sottospazio vettoriale.
\end{prop}
\begin{prop}
    Sia $(V, \mathbb{K}, +, \cdot)$ uno spazio vettoriale sul campo $\mathbb{K}$. Se $A \subseteq V$ un sottospazio vettoriale di $V$, allora
    $$
        \text{Span}(A) = A.
    $$
\end{prop}
\begin{prop}
    Sia $(V, \mathbb{K}, +, \cdot)$ uno spazio vettoriale sul campo $\mathbb{K}$. Siano $A, B \subseteq V$ insiemi, allora
    $$
        A \subseteq \text{Span}(B) \implies \text{Span}(A) \subseteq \text{Span}(B).
    $$
\end{prop}
\begin{prop}
    Sia $(V, \mathbb{K}, +, \cdot)$ uno spazio vettoriale sul campo $\mathbb{K}$. Siano $A, B \subseteq V$ insiemi, allora
    $$
        A \subseteq B \implies \text{Span}(A) \subseteq \text{Span}(B).
    $$
\end{prop}
\begin{exercise}
    \label{ex:1.3} Mostrare le quattro proposizioni qua sopra.
\end{exercise}
\section{Vettori linearmente indipedenti}
Un concetto fondamentale che ci aiuterà a capire meglio gli spazi vettoriali è il concetto di \textbf{vettori linearmente indipendenti}, che possiamo definire come segue:
\begin{definition}[vettori linearmente indipendenti]
    Sia $(V, \mathbb{K}, +, \cdot)$ uno spazio vettoriale sul campo $\mathbb{K}$ e siano $\underline{v_1}, \ldots, \underline{v_n} \in V$. Diremo che i vettori $\underline{v_1}, \ldots, \underline{v_n}$ sono linearmente indipendenti se l'equazione
    $$
    \alpha_1 \underline{v_1} + \ldots + \alpha_n \underline{v_n} = \underline{0}
    $$
    ha come \textbf{unica} soluzione quella banale, ovvero $\forall \, i \in \{ 1, \ldots, n \}, \alpha_i = 0$
\end{definition}
Perché siamo interessati a lavorare con i vettori linearmente indipendenti? La motivazione è semplice: le coordinate rispetto a tali vettori sono univoche, come mostriamo adesso
\begin{prop}
    Sia $(V, \mathbb{K}, +, \cdot)$ uno spazio vettoriale sul campo $\mathbb{K}$ e siano $\underline{v_1}, \ldots, \underline{v_n}$ dei vettori linearmente indipedenti, allora se si ha che
    $$
        \alpha_1 \underline{v_1} + \ldots + \alpha_n \underline{v_n} = \alpha_1' \underline{v_1} + \ldots + \alpha_n' \underline{v_n} \implies \forall i \in \{ 1,\ldots, n \}, \, \alpha_i = \alpha_i'
    $$
\end{prop}
\begin{proof}
    Portando tutti i termini a sinistra otteniamo una relazione equivalente,
    $$
        \alpha_1 \underline{v_1} + \ldots + \alpha_n \underline{v_n} = \alpha_1' \underline{v_1} + \ldots + \alpha_n' \underline{v_n} \iff (\alpha_1 - \alpha_1')\underline{v_1} + \ldots + (\alpha_n - \alpha_n') \underline{v_n} = \underline{0},
    $$
    ma, per ipotesi, abbiamo che i vettori $\underline{v_1}, \ldots, \underline{v_n}$ sono linearmente indipendenti, dunque l'unica soluzione a quest'ultima equazione è quella banale, ovvero deve risultare che
    $$
        \forall i \in \{ 1, \ldots, n \}, \alpha_i - \alpha_i' = 0 \implies \forall i \in \{ 1, \ldots, n \}, \alpha_i = \alpha_i',
    $$
    ovvero la tesi.
\end{proof}
\begin{remark}
    Essere vettori linearmente indipedenti \textbf{non} è una proprietà dei singoli vettori, ma una proprietà che hanno i vettori come insieme. Consiglio di guardare il prossimo esempio per chiarire meglio questo concetto
\end{remark}
\begin{example}[$\mathbb{R}^3$ "chiarificatore" di tutto il creato]
    Fissiamo $V = \mathbb{R}^3$ e consideriamo i vettori $\underline{v_1}$
\end{example}
	\clearpage
	\cleardoublepage
	\pagestyle{plain}
	\thispagestyle{empty}
	\vspace*{\fill}
	\begin{figure}[H]
		\centering
		\hspace*{-1cm}
		%\includegraphics[scale=0.25]{immagini/gamma_illustrata.png}
	\end{figure}
	\vspace*{\fill}
	\pagestyle{empty}
	\begin{appendices}
		\pagestyle{fancy}
		\fancyhead[LO]{\textbf{\thepage}}  % Numero di pagina in grassetto a sinistra
		\fancyhead[RO]{\nouppercase\rightmark}         % Capitolo e sezione a destra

		% Definizione dell'intestazione per pagine pari (Even)
		\fancyhead[LE]{\nouppercase\leftmark}          % Capitolo e sezione a sinistra
		\fancyhead[RE]{\textbf{\thepage}}  % Numero di pagina in grassetto a destra

		% Cancella il footer
		\fancyfoot{}
		\setlength{\headheight}{14.49998pt}
		\addtolength{\topmargin}{-2.49998pt}
		% Personalizzazione della linea orizzontale
		\renewcommand{\headrulewidth}{0.4pt}
		\renewcommand{\footrulewidth}{0pt}
	\end{appendices}
	\cleardoublepage
	\addcontentsline{toc}{chapter}{Bibliografia}
	\bibliographystyle{unsrt}
	\bibliography{bibliography.bib}
\end{document}